% Methods — Phases 1 & 2 only (Projet7_brainTumorDetection.ipynb)
\section{Methods}
\label{sec:methods}

\subsection{Dataset and preprocessing}
\label{sec:dataset}
We use a publicly available brain-tumour MRI collection (Kaggle Brain Tumor Detection \& Classification) with three classes---glioma, meningioma, and pituitary tumour---balanced at 500 slices per class (1{,}500 images total).
Images are converted to grayscale, resized to $128\times128$, cropped to remove black borders ($\text{thr}=0.05$), and contrast-enhanced with CLAHE before feature extraction.
The dataset is stratified into 80\%/20\% train/test (fixed random seed); 5-fold stratified cross-validation on the training fold selects hyperparameters.
Class balance is verified in Fig.~\ref{fig:class_balance} (imbalance ratio $\approx 1.0$; SMOTE not applied).

\begin{figure}[!t]
\centering
\includegraphics[width=0.55\linewidth]{class_balance_train}
\caption{Training-set class distribution after stratified split (500 samples per class).}
\label{fig:class_balance}
\end{figure}

\subsection{Separability metrics for classifier-free selection}
\label{sec:separability}
Phase~1 ranks descriptor configurations \emph{before} any classifier is trained.
For each hyperparameter setting we extract features on the training set and compute three complementary metrics:
\begin{itemize}
    \item \textbf{Fisher's discriminant ratio (FDR):} $\text{FDR}=\frac{1}{p}\sum_{j=1}^{p}\frac{\sigma_B^2(j)}{\sigma_W^2(j)+\varepsilon}$, measuring between-class versus within-class variance per feature dimension.
    \item \textbf{Mutual information (MI):} $\text{MI}(X_j;Y)=\sum p(x,y)\log\frac{p(x,y)}{p(x)p(y)}$, capturing non-linear dependency between each feature and the class label.
    \item \textbf{Davies--Bouldin index (DBI):} cluster compactness relative to inter-cluster separation (lower is better).
\end{itemize}
Each metric is min--max normalised to $[0,1]$ on the configuration grid; DBI is inverted so that higher is better.
The composite Phase~1 score is $(\text{FDR}_n + \text{MI}_n + (1-\text{DBI}_n))/3$.
The highest-scoring setting per family is locked for Phase~2 (Table~\ref{tab:phase1}); the full search grid is in Table~\ref{tab:descriptor_grids}.

% Hyperparameter search grids — Phase 1 (Projet7_brainTumorDetection.ipynb)
\begin{table*}[!t]
\centering
\caption{Phase~1 hyperparameter search space per descriptor family ($213$ configurations total). Bold values denote the selected optimum (Table~\ref{tab:phase1}).}
\label{tab:descriptor_grids}
\small
\begin{tabular}{@{}p{1.1cm}p{2.2cm}p{5.8cm}cp{2.8cm}@{}}
\toprule
Family & Hyperparameter & Search values & \# & Optimal \\
\midrule
\multirow{4}{*}{GLCM}
 & \texttt{distances} & \{1\}, \{1,2\}, \{1,2,4\}, \textbf{\{1,2,4,8\}}, \{2,4\}, \{4,8\} & 6 & \textbf{[1,2,4,8]} \\
 & \texttt{angles} (rad) & \{0\}, \{0,$\pi/2$\}, \{0,$\pi/4$,$\pi/2$,$3\pi/4$\}, 6-angle set & 4 & \textbf{[0]} \\
 & \texttt{levels} & 64, \textbf{128}, 256 & 3 & \textbf{128} \\
 & \texttt{symmetric} & \textbf{True}, False & 2 & \textbf{True} \\
\midrule
\multirow{3}{*}{LBP}
 & \texttt{P} (points) & \textbf{8}, 12, 16, 24 & 4 & \textbf{8} \\
 & \texttt{R} (radius) & \textbf{1}, 2, 3 & 3 & \textbf{1} \\
 & \texttt{method} & \textbf{uniform}, \texttt{nri\_uniform}, \texttt{ror} & 3 & \textbf{uniform} \\
\midrule
\multirow{2}{*}{DWT}
 & \texttt{wavelet} & haar, db1, db2, db4, sym2, coif1, \textbf{bior1.3} & 7 & \textbf{bior1.3} \\
 & \texttt{level} & \textbf{1}, 2, 3 & 3 & \textbf{1} \\
\midrule
\multirow{4}{*}{HOG}
 & \texttt{orientations} & \textbf{6}, 9, 12 & 3 & \textbf{6} \\
 & \texttt{pixels\_per\_cell} & (8,8), \textbf{(16,16)} & 2 & \textbf{(16,16)} \\
 & \texttt{cells\_per\_block} & \textbf{(2,2)}, (3,3) & 2 & \textbf{(2,2)} \\
 & \texttt{block\_norm} & L2-Hys & 1 & L2-Hys \\
\bottomrule
\end{tabular}
\end{table*}


\subsection{Grey-level co-occurrence matrix (GLCM)}
\label{sec:glcm}
The GLCM~\cite{ref_fonc20231248452,ref_hsr271195} tabulates how often grey levels $i$ and $j$ co-occur at a given pixel offset $(d,\theta)$.
For an image quantised to $L$ levels, the matrix $\mathbf{P}$ has size $L\times L$; we average over all distance--angle pairs and extract six Haralick properties: contrast, dissimilarity, homogeneity, energy, correlation, and angular second moment (ASM).
\textbf{Output:} 6 features.
\textbf{Optimal:} $d\in\{1,2,4,8\}$, $\theta=0$, $L=128$, symmetric=True (score $=0.781$).

\subsection{Local binary patterns (LBP)}
\label{sec:lbp}
LBP~\cite{ref_fnins2021634926} encodes local texture by thresholding a circular neighbourhood of $P$ pixels on a ring of radius $R$.
\textbf{Output:} $P{+}2$ histogram bins for \texttt{uniform} method (10 dims for $P=8$).
\textbf{Optimal:} $P=8$, $R=1$, \texttt{uniform} (score $=1.000$).

\subsection{Discrete wavelet transform (DWT)}
\label{sec:dwt}
The 2-D DWT decomposes the image into approximation (LL) and detail sub-bands (LH, HL, HH) at each scale~\cite{ref_2409,ref_iasc2022024635}.
We compute mean, standard deviation, energy, and Shannon entropy per sub-band.
\textbf{Output:} $4\times(3L{+}1)$ features (16 at $L=1$).
\textbf{Optimal:} \texttt{bior1.3}, level $=1$ (score $=0.994$).

\subsection{Histogram of oriented gradients (HOG)}
\label{sec:hog}
HOG~\cite{ref_s41598024772437,ref_wo2023124612} accumulates gradient-orientation histograms per cell with block normalisation.
\textbf{Output:} 1{,}176 dimensions with optimal settings on $128\times128$ inputs.
\textbf{Optimal:} 6 orientations, $(16,16)$ cells, $(2,2)$ blocks (score $=0.997$).

Fig.~\ref{fig:descriptor_demo} illustrates all four descriptors on one preprocessed slice using these optimal parameters.

\begin{figure*}[!t]
\centering
\includegraphics[width=\linewidth]{descriptor_demo}
\caption{Visualisation of Phase~1 optimal descriptors on one MRI slice: preprocessed input, GLCM co-occurrence matrix with Haralick properties, LBP code map, DWT level-1 sub-bands, and HOG energy overlay.}
\label{fig:descriptor_demo}
\end{figure*}

% Auto-generated by tables_to_latex.py — do not edit by hand
\begin{table}[!t]
\centering
\caption{Phase 1 optimal feature extractor parameters.}
\label{tab:phase1}
\begin{tabular}{@{}llcc@{}}
\toprule
Extractor & Best config & Dim & Score \\
\midrule
GLCM & d=[1, 2, 4, 8] a=1ang l=128 sym=True & 6 & 0.7809 \\
LBP & P= 8 R=1 method=uniform & 10 & 1.0000 \\
DWT & wavelet=bior1.3  level=1 & 16 & 0.9941 \\
HOG & ori=6 ppc=(16, 16) cpb=(2, 2) norm=L2-Hys & 1176 & 0.9969 \\
\bottomrule
\end{tabular}
\end{table}


\subsection{Feature vectors and cleaning}
\label{sec:features}
Eight feature sets are built from Phase~1 optimal extractors plus first-order statistics:
A (Stats), B (Tex/opt = GLCM+LBP), C (DWT/opt), D (HOG), and unions A+B, B+C, A+B+C, Full(opt) ($\sim$1{,}987 raw dimensions).
Before training: \texttt{VarianceThreshold}, Pearson correlation filter at 0.95, \texttt{StandardScaler}.

\subsection{Phase 2: Classifier benchmark}
\label{sec:phase2}
Eight classifiers with \texttt{GridSearchCV} (macro-F1, 5-fold CV): SVM, KNN, RF, XGBoost, LightGBM, LR, MLP, ExtraTrees.
% Auto-generated by tables_to_latex.py — do not edit by hand
\begin{table}[!t]
\centering
\caption{Phase 2 classifier hyperparameter search spaces (per model, before $\textbackslash{}times$ eight feature sets).}
\label{tab:phase2_models}
\begin{tabular}{@{}llcc@{}}
\toprule
Model & Key hyperparameters & Candidates & $\times$ 5-fold CV \\
\midrule
SVM & 4 kernels, C, $\gamma$, degree, coef$_0$ & 903 & 4,515 \\
KNN & n\_neighbors, metric, weights, p & 384 & 1,920 \\
Random Forest & n\_estimators, max\_depth, max\_features, criterion & 14,112 & 70,560 \\
XGBoost & n\_estimators, max\_depth, lr, subsample, colsample & 6,480 & 32,400 \\
LightGBM & n\_estimators, max\_depth, lr, num\_leaves, subsample & 13,824 & 69,120 \\
Logistic Reg. & C, penalty (L1/L2/elastic-net), class\_weight & 108 & 540 \\
MLP & hidden\_layers, activation, $\alpha$, lr, batch & 432 & 2,160 \\
ExtraTrees & n\_estimators, max\_features, min\_samples & 162 & 810 \\
\midrule
\textbf{Total} & & \textbf{36,405} & \textbf{182,025} \\
\bottomrule
\end{tabular}
\end{table}


\subsection{Evaluation and statistics}
Primary metric: macro-averaged F1.
We report Cohen's $\kappa$, bootstrap 95\% CIs, McNemar, Wilcoxon, and Friedman--Nemenyi (Section~\ref{sec:statistics}).
t-SNE (Fig.~\ref{fig:tsne}) visualises class separability in Full(opt).

\begin{figure}[!t]
\centering
\includegraphics[width=\linewidth]{tsne_best_feature_set}
\caption{t-SNE projection of the Full(opt) training features (subsampled), coloured by tumour class.}
\label{fig:tsne}
\end{figure}
